\documentclass[10pt,a4paper]{amsart}
\usepackage[margin=19mm]{geometry}
\usepackage{amsmath,amssymb,mathtools}
\usepackage[scr=rsfs,cal=euler]{mathalfa}
\usepackage{xfrac}
\usepackage[unicode,hidelinks,bookmarks=false]{hyperref}
\newcommand{\ie}{i.e.\ }
\newcommand{\cf}{cf.\ }
\newcommand{\resp}{respectively\ }
\newcommand{\Subsection}{\subsection}
\providecommand{\nobreakdash}{\nobreak\hbox{-}}
\setlength{\emergencystretch}{2em}
\multlinegap0pt
\usepackage{bm}
\usepackage{bbm}
\usepackage{dsfont}
\usepackage{xspace}
\usepackage{cancel}
\usepackage{mathtools}
\usepackage{amsthm}
\makeatletter
\@ifundefined{theorem}{\newtheorem{theorem}{Theorem}}{}
\@ifundefined{lemma}{\newtheorem{lemma}[theorem]{Lemma}}{}
\makeatother
\makeatletter
\newcommand{\Om}{\Omega}
\newcommand{\R}{{\mathbb R}}
\renewcommand{\a}{\alpha}
\newcommand{\curl}{\operatorname{curl}}
\renewcommand{\div}{\operatorname{div}} % {\mathrm{div}\,}
\expandafter\ifx\csname ex\endcsname\relax
\newenvironment{ex}{\pushQED{\qed} \exbase}{\popQED\endexbase}
\fi
\newcommand{\grad}{\nabla}
\newcommand{\la}{\langle}
\newcommand{\lm}{\lambda}
\newcommand{\pa}{\partial}
\newcommand{\ra}{\rangle}
\expandafter\ifx\csname remark\endcsname\relax
\newenvironment{remark}{\pushQED{\qed} \remarkbase}{\popQED\endremarkbase}
\fi
\makeatother
\title[A rigidity result for the 3D capillary liquid drop with cons]{A rigidity result for the 3D capillary liquid drop with constant vorticity}
\author{Pietro Baldi, Domenico Angelo La Manna, Giuseppe La Scala}
\date{Source: arXiv:2607.00450v1 (CC BY 4.0)}
\begin{document}
\maketitle

\noindent\emph{Source and licence.} A rigidity result for the 3D capillary liquid drop with constant vorticity. Version arXiv:2607.00450v1 (CC BY 4.0), distributed under the Creative Commons Attribution 4.0 International licence, as recorded in the arXiv licence metadata for this exact version and shown on the abstract page https://arxiv.org/abs/2607.00450v1. Full rights notice: licenses/arxiv-2607.00450-CC-BY-4.0.txt.\par\smallskip

\noindent\textit{Editorial scope.} Counted unit: the Lemma whose source label is \texttt{lem:beta} in the article (source0.tex lines 1989--1996, source label \texttt{lem:beta}), delivered together with the article's own complete original proof of it (source0.tex lines 1998--2130, one \texttt{proof} environment of 133 lines). No mathematics is written, repaired or re-derived: statement and proof are the article's own text line for line. Required context retained uncounted: time.indep.sol (lines 463--465); steady.rotating.solution (lines 541--553); thm:global (lines 557--583); eq:tang.divergence (lines 911--914); torsion.gradient (lines 965--967); torsion.integral (lines 969--979); thm:4.1 (lines 1193--1222); eq:torsion (lines 1206--1211); eq:over.determined (lines 1213--1219); eq:defn.beta (lines 1827--1829); eq:constant.c (lines 1940--1942). Every definition, display and label of those lines is retained unchanged. Omitted from this package and not redistributed: 1--462, 466--540, 554--556, 584--910, 915--964, 968--968, 980--1192, 1212--1212, 1220--1826, 1830--1939, 1943--1988, 1997--1997, 2131--2705. Nothing inside a retained excerpt is omitted or altered: every retained body is delivered exactly as the article's lines are written, so no omission receipt of any kind accompanies this package. Citations: the retained excerpts carry no citation command at all, so no bibliography entry of the article is transcribed and no reference is redistributed. Reference adaptation: none was needed and none was made; every reference inside the delivered unit resolves inside it, so no unresolved reference or citation is produced. Printed slips kept literal: no printed word, symbol, hypothesis or number of the counted statement or of its proof was altered. Layout and class: the article's own class is \texttt{article} with packages including anysize, verbatim, babel, amsmath, amsfonts, amssymb, amsthm, xcolor, comment, hyperref, mathrsfs, dsfont, accents; the wrapper is the lane's pinned \texttt{amsart} editorial wrapper at 10pt on A4, which restates the theorem-like environments the retained text needs and adds the article's own macro definitions that the retained text needs (\texttt{\textbackslash Om}, \texttt{\textbackslash R}, \texttt{\textbackslash a}, \texttt{\textbackslash curl}, \texttt{\textbackslash div}, \texttt{\textbackslash grad}, \texttt{\textbackslash la}, \texttt{\textbackslash lm}, \texttt{\textbackslash pa}, \texttt{\textbackslash ra}), with extra wrapper packages (bm, bbm, dsfont, xspace, cancel, mathtools). The delivered numbering is the unit's own. Assets: no retained span contains a figure environment, a graphics-inclusion command, a PDF-page inclusion, a table or a TikZ picture; no image file is copied into this package, the package declares no asset dependency and no image file is redistributed with it. Licence: version arXiv:2607.00450v1 of this article is distributed under the Creative Commons Attribution 4.0 International licence, as recorded in the arXiv licence metadata for this exact version and shown on the abstract page https://arxiv.org/abs/2607.00450v1. A full rights notice is delivered at licenses/arxiv-2607.00450-CC-BY-4.0.txt. Characters: every retained excerpt is pure ASCII, so the delivered engine, XeLaTeX with its default Latin Modern text font, reports no missing character.
\begin{equation} \label{time.indep.sol}
\Om, u = \text{independent of time}.
\end{equation}

\begin{equation} \label{steady.rotating.solution}
\begin{cases}
\la u, \nabla \ra  u=-\nabla p   & \text{in } \Omega,
\\
\div u=0 & \text{in } \Omega,
\\
\curl u= \alpha_0 e_3 & \text{in } \Omega,
\\
p=\sigma_0 H_{\pa \Omega} &\text{on } \pa \Omega,
\\
\la u ,\nu_{\pa\Om}\ra =0  &\text{on } \pa \Omega.
\end{cases}
\end{equation}

\begin{theorem}\label{thm:global}
Let $\Omega$ be a strictly convex open set with $C^2$ boundary 
and $u\in C^2(\Omega,\R^3) \cap C(\overline{\Om}, \R^3)$. 
Suppose that $(\Omega,u)$ is a solution of \eqref{time.indep.sol}, \eqref{steady.rotating.solution}.
Then 
\begin{itemize}

\item[$(i)$]
There exists a real constant $\lambda_0$, a strictly convex open set $D\subset \R^2$ 
and a real function $f \in C^2(D) \cap C(\overline{D})$, 
with $f>0$ in $D$ and $f=0$ on $\pa D$, such that
\[
\Omega = \{(x', x_3) \in \R^3 : x' \in D, \ |x_3-\lambda_0| < f(x') \}.
\]
Moreover the vector field $u$ is independent of $x_3$, 
and it has third component $u_3 = 0$. 

\item[$(ii)$]
If 
\begin{equation}\label{eq:sigma.alpha}
\frac{\alpha_0^2}{\sigma_0} |D|^{\frac32} \leq \sqrt{2\pi^3}, 
\end{equation}
then $D$ is a disk, 
$\pa\Om$ is a surface of revolution, 
and the velocity vector field is $u = \frac{\a_0}{2} (-x_2, x_1, 0)$. 
\end{itemize}
\end{theorem}

\begin{equation}\label{eq:tang.divergence}
\int_{\pa E} \div_{\pa E} X \, d\sigma
= \int_{\pa E} H_{\pa E} \la X, \nu_{\pa E} \ra \, d\sigma
\end{equation}

\begin{equation}\label{torsion.gradient}
\nabla v_E(x) = -|\nabla v_E(x)| \nu_{\pa E}(x).
\end{equation}

\begin{align}
& \int_E v_E \, dx = \int_E v_E (-\Delta v_E) \, dx 
= \int_E |\grad v_E|^2 \, dx, 
\label{torsion.integral} 
\\ 
& \int_{\pa E} |\nabla v_E| \, d\sigma 
= \int_{\pa E} - \la \grad v_E, \nu_{\pa E} \ra \, d\sigma 
= - \int_E \Delta v_E \, dx 
= |E|. 
\label{torsion.gradient.integral}
\end{align}

\begin{theorem}[An overdetermined problem for the torsion function]
\label{thm:4.1}
Let $(\Omega,u)$ satisfy the assumptions of Theorem \ref{thm:global}, 
and suppose that the barycenter of $\Omega$ is at the origin. 
Let $\lm_0, D, f$ be given by item $(i)$ of Theorem \ref{thm:global}, 
so that $\lm_0 = 0$, 
\begin{equation} \label{eq:omega.is.graph}
D = \{ x' \in \R^2 : (x',0) \in \Omega \}, 
\quad \ 
\pa \Omega \cap \{x_3>0\} = \{(x', f(x')) : x' \in D\}.
\end{equation}
Then the torsion function $v_D \in H^1_0(D)$ of the planar set $D$, 
that is, the solution of the problem 
\begin{equation} \label{eq:torsion} 
\begin{cases}
\Delta v_D = -1 & \text{in } D, \\
v_D=0 & \text{on } \pa D,
\end{cases}
\end{equation}
also satisfies the equation 
\begin{equation} \label{eq:over.determined}
\frac{\alpha_0^2}{2}  | \nabla v_D|^{2}  
- \frac{ \sigma_0 \beta }{ | \nabla v_D| }  
+ \sigma_0 H_{\pa D}
= c 
\quad \ \text{on } \pa D,
\end{equation}
where $H_{\pa D}$ is the curvature of the boundary $\pa D$ (which is a planar curve),
and $\beta$ and $c$ are real constants depending only on $D$.
\end{theorem}

\begin{equation}\label{eq:defn.beta}
\beta := \lim_{s \to 0^+} \frac{ \tilde f''(s) }{ (\tilde f'(s))^3 } \in \R.
\end{equation}

\begin{equation}\label{eq:constant.c}
c=\frac{1}{|D|} \left(\frac{3\alpha^2_0}{2} \int_{D} |\nabla v_D|^2\,dx +\sigma_0 P(D)  \right).
\end{equation}

\begin{lemma}\label{lem:beta}
Assume the hypotheses of Theorem \ref{thm:4.1}. 
The constant $\beta$ defined in \eqref{eq:defn.beta} is given by 
\begin{equation} \label{formula.in.lem:beta}
\beta = - \Big( \sigma_0 \int_{\pa D} \frac{\la x,\nu_{\pa D}\ra}{|\nabla v_D|} \, d\sigma \Big)^{-1} 
\Big\{ \sigma_0 P(D) + \alpha_0^2 \int_D |\nabla v_D|^2 \, dx \Big\}.
\end{equation}
\end{lemma}

\begin{proof} 
Multiplying \eqref{eq:over.determined} by $\la x,\nu_{\pa D}\ra$ 
and integrating along the curve $\pa D$, we get 
\begin{equation} \label{along.01}
\frac{\alpha_0^2}{2} \int_{\pa D} | \nabla v_D|^{2} \la x, \nu \ra \, d\sigma 
- \sigma_0 \beta \int_{\pa D} \frac{ \la x, \nu \ra }{ | \nabla v_D| } \, d\sigma
+ \sigma_0 \int_{\pa D} H_{\pa D} \la x, \nu \ra \, d\sigma
= c \int_{\pa D} \la x, \nu \ra \, d\sigma.
\end{equation}
We calculate each term. 
%to find
%\[
%\frac{\alpha_0^2}{2}\int_{\pa D} \la x,\nu_{\pa D} \ra |\nabla v_D |^2 -\beta \sigma_0\int_{\pa D}\frac{ \la x,\nu_{\pa D} \ra }{|\nabla v_D |}\, d\sigma 
%+\sigma_0 P(D)=2c|D|
%\]
%Recall the general identity $\div (x \ph) 
%= (\div x) \ph + \la x, \grad \ph \ra 
%= 2 \ph + \la x, \grad \ph \ra$, which holds for any scalar function $\ph$ of two variables. 
By the divergence theorem,
\begin{align*}
\int_{\pa D} |\nabla v_D|^2 \la x,\nu_{\pa D} \ra \, d\sigma
& = \int_D \div \{ x |\nabla v_D|^2 \} \, dx 
\\ 
& = \int_D (\div x) |\nabla v_D|^2\, dx 
+ \int_D \la x, \nabla ( |\nabla v_D|^2 ) \ra \, dx
\\ 
& = 2 \int_D |\nabla v_D |^2 \, dx 
+ 2 \int_D \la (\nabla^2 v_D) (\nabla v_D) , x \ra \, dx,
\end{align*}
where $\grad^2 v_D$ is the Hessian matrix of $v_D$. 
From the general identity
\[
\div \{ \la \grad w , x \ra \grad w \} 
= \la \grad w , x \ra \Delta w 
+ |\grad w|^2 
+ \la (\grad^2 w)(\grad w), x \ra
\]
applied to the torsion function $v_D$, using the identity $\Delta w_D = -1$, we find 
\[
\int_D \la (\nabla^2 v_D) (\nabla v_D), x\ra\,dx
= \int_D \div \{ \la \nabla v_D, x \ra \nabla v_D \} \, dx 
+ \int_D \la \nabla v_D, x \ra \, dx 
- \int_D |\nabla v_D|^2 \, dx.
\]
By the divergence theorem and formula \eqref{torsion.gradient} for the gradient of $v_D$,  
\[
\int_D \div \{ \la \nabla v_D, x \ra \nabla v_D \} \, dx 
= \int_{\pa D} \la \nabla v_D, x \ra \la \nabla v_D , \nu_{\pa D} \ra \, d\sigma
= \int_{\pa D} |\nabla v_D|^2 \la x,\nu_{\pa D}\ra \, d\sigma.
\]
By the divergence theorem applied to the vector field $x v_D$, 
which has divergence $\div(x v_D) = (\div x) v_D + \la x, \grad v_D \ra 
= 2 v_D + \la x, \grad v_D \ra$ in $D$ 
and boundary value $x v_D = 0$ on $\pa D$ by \eqref{eq:torsion},    
one has 
\[
0 = \int_{\pa D} \la x v_D , \nu_{\pa D} \ra \, d\sigma 
= \int_D \div(x v_D) \, dx 
= 2 \int_D v_D \, dx + \int_D \la x , \grad v_D \ra \, dx,
\]
and, by \eqref{torsion.integral},
\[
\int_D \la x , \grad v_D \ra \, dx 
= - 2 \int_D v_D \, dx 
= - 2 \int_D |\grad v_D|^2 \, dx.
\]
From these identities, we obtain  
%using the divergence theorem, 
%\eqref{eq:torsion}, \eqref{torsion.gradient}, \eqref{torsion.integral},
%we find 
%\begin{align*}
%\int_D \la (\nabla^2 v_D) (\nabla v_D), x\ra\,dx
%& = \int_D \div \{ \la \nabla v_D, x \ra \nabla v_D \} \, dx 
%- \int_D \la \nabla v_D, x \ra \Delta v_D \, dx 
%- \int_D |\nabla v_D|^2 \, dx
%\\
%& = \int_{\pa D} \la \nabla v_D, x \ra \la \nabla v_D , \nu_{\pa D} \ra \, d\sigma
%+ \int_D \la \nabla v_D, x \ra \, dx  
%- \int_D |\nabla v_D|^2 \, dx
%\\
%& = \int_{\pa D} |\nabla v_D|^2 \la x,\nu_{\pa D}\ra \, d\sigma
%+ \int_D \la \nabla v_D, x \ra \, dx  
%- \int_D |\nabla v_D|^2 \, dx.
%\end{align*}
%Moreover, by using $\grad^2v_D=(\grad^2v_D)^T$, 
%divergence theorem, 
%\eqref{eq:torsion}, \eqref{torsion.gradient} and \eqref{torsion.integral},
%\[
%\begin{split}
 %\int_D \la \nabla^2 v_D \nabla v_D, x\ra\,dx
%&=\int_D \la (\nabla^2v_D)^T  \nabla v_D, x\ra\,dx
%\\&=\int_{D} \Big( \div \{ \la\nabla v_D,x\ra \nabla v_D \} - (\Delta v_D) \la\nabla v_D,x\ra 
%- |\nabla v_D|^2 \Big)\, dx
%\\
%&=\int_{\pa D}  \la\nabla v_D,x\ra \la \nabla v_D,\nu_{\pa D} \ra\, d\sigma
%+\int_{D}(\la \nabla v_D,x\ra - |\nabla v_D|^2) dx\\
%&=\int_{\pa D} |\nabla v_D|^2 \la x,\nu_{\pa D}\ra d\sigma
%+\int_{D}( -2v_D\, -|\nabla v_D|^2) dx
%\\&
%=\int_{\pa D} |\nabla v_D|^2 \la x,\nu_{\pa D}\ra d\sigma
%-3\int_{D}|\nabla v_D|^2 dx.
%\end{split}
%\]
\begin{equation}\label{eq:boundary.interior}
\int_{\pa D}|\nabla v_D|^2\la x,\nu_{\pa D}\ra \, d\sigma 
= 4 \int_{D}|\nabla v_D|^2\,dx.
\end{equation}
By \eqref{eq:tang.divergence} with $X(x)=x$, 
since $\div_{\pa D} x=1$, we get
\[
\int_{\pa D} H_{\pa D} \la x, \nu_{\pa D} \ra \, d\sigma = P(D), 
\]
and, from the divergence theorem, 
\[
\int_{\pa D} \la x, \nu_{\pa D} \ra \, d\sigma 
= \int_D \div x \, dx 
= \int_D 2 \, dx 
= 2 |D|.
\]
Hence \eqref{along.01} becomes 
\[
2 \alpha_0^2 \int_{D}|\nabla v_D|^2\,dx 
- \sigma_0 \beta \int_{\pa D} \frac{\la x, \nu_{\pa D} \ra }{ |\grad v_D| } \, d\sigma 
+ \sigma_0 P(D) = 2 c |D|.
\]
Replacing the constant $c$ with its formula in \eqref{eq:constant.c}, 
we get the result. 
%\[
%\frac{\alpha_0^2}{4}\int_{\pa D}  |\nabla v_D|^2\la x,\nu_{\pa D}\ra\,d\s + \beta \sigma_0 \int_{\pa D}\frac{ \la x,\nu_{\pa D} \ra }{|\nabla v_D|}\, d\sigma
%+\sigma_0 P(D)=0
%\]
%which readily gives the result.
\end{proof}

\end{document}
